\section{The triangular induction and its determinant interface}
\label{sec:triangle}

\subsection{Keep the mixed Gram identity}

Let \(V\) be the ambient space of principal forms and
\(s:U\hookrightarrow V\) the target subspace, with \(\dim U=d_T\).
Let \(j_i\) identify \(V\) with the bottom of row \(i\), and let
\(X_{i,n}\) be a weighted row of \(n\) lifts. Set
\[
 q_n(\vartheta)=
 \frac{w_0(\vartheta)w_1(\vartheta)^n}{w_n(\vartheta)}.
\]
The induction should retain a statement of the form
\begin{equation}
 \ip{X_{i,n}j_i x}{X_{i,n}j_i(sy)}
   =q_n\,\ip{x}{s(K_{i,n}y)}
 \qquad (x\in V,\ y\in U),
 \label{eq:mixed}
\end{equation}
with the exact weights of the existing triangle and an endomorphism witness
\(K_{i,n}\) of \(U\). Dependence on the twist is suppressed here.

Quantifying \(x\) over all of \(V\) is essential. In the induction, the
lift--descent composite carries \(x\) to another ambient vector \(x'\).
A Gram identity only on \(U\times U\) would not suffice without a new
invariance theorem. The source already supplies the stronger identity:
see \code{IsTriangle.exists_inner_weightedChain_eq} in
\src{New/Determinant/Chain/TriangleNorms}{\code{TriangleNorms.lean}}.

The one-step mixed identity has the form
\begin{equation}
 w_0\ip{a_i j_i x}{a_i j_i(sy)}
       =C_i\ip{x}{s(S_i y)}.
\end{equation}
Use the existing square-commutation and square-adjoint laws to move a
descent across the row, apply \eqref{eq:mixed} to the shorter row, and then
apply (6.2). The weight quotients cancel to \(q_{n+1}\). This gives
\begin{equation}
 K_{i,0}=I,\qquad
 K_{i,n+1}=a_{i,n}S_iK_{i+1,n},
 \qquad a_{i,n}=\kappa_{i,n}\beta_{i,n}C_i.
\end{equation}
The constants \(\beta,\kappa\) are the accumulated commutation and adjoint
scalars. There is no need for the public theorem to identify \(K_{i,n}\)
with a named product of all \(S_i\).

\begin{proposition}[Determinants need no operator commutativity]
\label{prop:det}
For square matrices of a common dimension \(d_T\) over a commutative ring,
the recurrence (6.3) gives
\begin{equation}
 \det K_{i,n}
   =A_{i,n}^{d_T}\prod_{j=i}^{i+n-1}\det S_j,
 \qquad
 A_{i,0}=1,\quad A_{i,n+1}=a_{i,n}A_{i+1,n}.
\end{equation}
The matrices \(S_j\) need not commute.
\end{proposition}
\begin{proof}
Apply \(\det(aT)=a^{d_T}\det T\) and
\(\det(ST)=\det S\,\det T\) at each induction step.
Only the scalar coefficient ring is commutative.
\end{proof}

This is formalized in \code{MatrixChainBudget.determinant_chain}.
The file does not purport to prove (6.3) from theta-lift geometry.
Equation \eqref{eq:mixed} specifies the source-level induction needed to
instantiate the checked algebraic lemma.

\subsection{Arbitrary budgets and the exact crossing count}

Suppose crossing \((r,k)\) has squared cost \(B_{rk}\).
The row of length \(n+1\) at \(i\) crosses the \(n\) squares
\((i+1,k)\), \(i+2\le k\le i+n+1\), and then uses the shorter row.
Its row cost is therefore
\[
 e_{i,n}=\sum_{k=i+2}^{i+n+1}B_{i+1,k},\qquad
 T_{i,0}=0,\quad T_{i,n+1}=e_{i,n}+T_{i+1,n}.
\]
The scalar accumulated over the full triangle has cost
\begin{equation}
 T_{0,m}=\sum_{1\le r<k\le m}B_{rk}.
 \label{eq:weightedcost}
\end{equation}
For uniform crossing budget \(B\),
\[
 T_{n+1}=nB+T_n,\qquad T_0=0,\qquad
 T_m=\binom m2 B.
\]
The archived Lean file proves the arbitrary double-sum identity, the uniform
recurrence, \(2\binom m2+m=m^2\), and propagation of guarded scalar
certificates. With \eqref{eq:onecrossing}, the final squared cost is
\begin{equation}
 2\binom m2 tx_1.
\end{equation}
Theorem~\ref{thm:root} gives the unsquared bound
\begin{equation}
 \HLe\!\left(\binom m2 tx_1,\ r_{\kappa,\theta}\right).
\end{equation}

The binomial count is already visible in the source's product over \(i<j\);
sharpening \(m^2\) to \(\binom m2\) is not the main novelty.
The simplification is that the induction no longer carries a separate
symbolic ratio product and global norm exponent to obtain this bound.
Keeping the individual \(B_{rk}\) in \eqref{eq:weightedcost} also
accommodates stronger local estimates without changing the induction.

\subsection{The exact boundary cancellation remains}

The triangle weights, finite group orders, and level indices are not
individually pure conductor powers. They contain Euler factors.
There is an exact calculation here that cannot be replaced by an assertion
that all normalization factors cost zero.

Write \(I(\vartheta)\) for the level index, \(M(\vartheta)\) for the
matrix-space cardinality, \(d_W\) for the fixed Weyl-index dimension,
\(g(\vartheta)\) for \code{gramFactor}, and
\(u(\vartheta)\) for \code{twistFactor}.
The original identities give
\begin{align}
 I(\vartheta)&=
  \gamma N(Q_\vartheta)^\ell g(\vartheta)^m
      \frac{w_m}{w_0w_1^m}, \label{eq:boundaryI}\\
 \frac{g(\vartheta)}{M(\vartheta)}&=g_0u(\vartheta),
 \label{eq:boundaryg}
\end{align}
with \(\gamma,g_0,\ell\) independent of the twist.
The first is \code{exists_stdLevelIndex_lastLevel_eq} in
\src{New/Determinant/Chain/ChainNorms}{\code{ChainNorms}}; the second is
\code{DescentPair.gramFactor_div_card} in
\src{New/Determinant/Norms/AdjointThetaLift}{\code{AdjointThetaLift}}.
The conversion from weighted-chain inner products to the actual chain
contributes \(I(\vartheta)d_W\).

In the application, the one-step operators in (6.3) include the factor
\(M(\vartheta)^{-1}\). Taking determinants gives
\begin{equation}
 Z(\vartheta)=
 \left[
 A_{0,m}(\vartheta)
 \frac{I(\vartheta)d_Ww_0w_1^m}
      {w_mM(\vartheta)^m}
 \right]^{d_T}
 \prod_{j=1}^m\det R_j(\vartheta).
 \label{eq:rawZ}
\end{equation}
The unchanged single-lift comparison is
\begin{equation}
 Z_j(\vartheta)=
   (A_j u(\vartheta))^{d_T}\det R_j(\vartheta),
 \qquad A_j\ne0,
\end{equation}
where \(A_j\) is twist-independent. Here \(Z_j(\vartheta)\) already denotes
the source's single lift evaluated at the appropriate \(j\)-th twist power.
All first \(m\) powers have the same conductor support.

Substitution of \eqref{eq:boundaryI}--\eqref{eq:boundaryg} into
\eqref{eq:rawZ} yields the boundary coefficient
\[
 \gamma d_W g_0^m\,N(Q_\vartheta)^\ell u(\vartheta)^m.
\]
The factor \(u(\vartheta)^m\) cancels exactly against (6.11).
Only the remaining pure norm power can then be discarded in the height
calculus. The equal-multiplicity comparison that moves each single-row
determinant to the common target remains necessary as well.

Taking absolute values of the positive Gram determinants now exports
\begin{equation}
 Z(\vartheta)=\rho(\vartheta)^{d_T}
       \prod_{j=1}^m Z_j(\vartheta),
 \qquad \rho(\vartheta)>0.
 \label{eq:Zfactor}
\end{equation}
After normalizing at the trivial twist, the inverse ratio
\(\rho(1)/\rho(\theta_\iota)\) is the normalized modulus of the accumulated
scalar times a pure conductor power. It has exactly the same height
allowance as (6.7). No global exponent or explicit
\code{quotientRatio} product needs to be exported.
The zero-target case is handled separately before asserting
\(\rho(\vartheta)>0\).

\section{Recovering Theorem C with unchanged constants}
\label{sec:C}

The same regular family \(\theta\) already satisfies, for every
\(1\le j\le m\),
\begin{equation}
 \HLe\!\left(
 d_T\left(\frac{L}{m}-k\,\operatorname{directionAmount}\right),
 \ \iota\mapsto\frac{Z_j(1)}{Z_j(\theta_\iota)}
 \right).
 \label{eq:singleheight}
\end{equation}
This is
\src{New/Determinant/Twist/LongForms}{\code{exists_hasSmallRatios_heightLe_liftDeterminant_div}}.
Adding the \(m\) contributions gives exactly \(d_T(L-H)\).

For the chain scalar, \eqref{eq:Zfactor} and (6.7) imply
\[
 \HLe\!\left(
 d_T\left(\binom m2 tx_1+L-H\right),
 \ \iota\mapsto \frac{Z(1)}{Z(\theta_\iota)}
 \right).
\]
The existing admissibility estimate \(m^2tx_1\le b\) gives the unchanged
public result
\begin{equation}
 \HLe\!\left(d_T(b+L-H),
 \ \iota\mapsto Z(1)/Z(\theta_\iota)\right).
 \label{eq:oldbound}
\end{equation}
There is no extra factor two and no rescaling of \(H\).

For auxiliary forms, keep the actual scalar
\(\gamma(\vartheta)\in\C^\times\), normalized by \(\gamma(1)=1\),
and their exact integrality identities. Replace only the explicit
\code{turnRatio} modulus formula by
\begin{equation}
 \HLe\!\left(
 2m^2tx_1,\ 
 \iota\mapsto(\norm{\gamma(\theta_\iota)}:\C)^{-2}
 \right).
\end{equation}
Here the consumer genuinely needs the squared reciprocal modulus, so
power descent is not required. The already used inequality
\(2m^2tx_1\le b\) recovers its original allowance \(b\).

Complete the independent family by \(n'\) auxiliary forms. The unchanged
archimedean bound is \(2x_1^2b^2\) per added form. With
\(n'\le4d_T/b^2\),
\[
 n'(2x_1^2b^2+b)
 \le d_T(8x_1^2+4/b)\le d_Tb.
\]
The existing product-formula theorem therefore gives
\begin{align*}
 0
 &\le V_0+2\delta+n'(2x_1^2b^2+b)+d_T(b+L-H)\\
 &\le V_0+2\delta+d_T(2b+L-H).
\end{align*}
Using \(b\le L\) proves \eqref{eq:theoremC}.
Every ordinary-direction contribution in \eqref{eq:singleheight} survives.
This is the concrete compatibility check that distinguishes a usable
height refactor from a formal rearrangement with a changed main term.

\section{Where the readable proof would become shorter}

\begin{center}
\begin{tabular}{@{}p{.33\textwidth}r p{.52\textwidth}@{}}
\toprule
Module (relative to \code{New/Determinant}) & Lines & Proposed change\\
\midrule
\code{Chain/Moduli} & 430 &
Replace the global modulus API by local squared-cost lemmas;
retain tame/primitive-character results and the arithmetic inputs.\\
\code{Chain/ChainNorms} & 892 &
Attach costs to the existing square and adjoint witnesses;
retain maps, exact normalizations, and injectivity inputs.\\
\code{Chain/TriangleNorms} & 364 &
Use a graded additive cost; expose mixed Gram witnesses and their
determinants instead of a symbolic global ratio product.\\
\code{Chain/NormOfChain} & 528 &
Retain one-lift determinant and boundary cancellation;
remove global-exponent extraction and ratio-product endpoints.\\
\code{AuxiliaryForms/Turns} & 1,002 &
Replace scalar modulus conclusions by squared costs; retain the
geometric turn and integrality proofs.\\
\code{AuxiliaryForms/DiagonalTurns} & 672 &
Propagate the costs with the same actual scalars.\\
\code{AuxiliaryForms/SpanningForms} & 796 &
Export a nonzero auxiliary scalar and its cost, keeping spanning
and size conclusions.\\
\code{Twist/LongForms} & 250 &
Consume two height certificates directly and retain Theorem C.\\
\bottomrule
\end{tabular}
\end{center}

These eight files contain 4,934 physical lines; that is the affected
\emph{neighborhood}, not a deletion estimate. Much of it proves geometry,
integrality, character compatibility, or nonvanishing and remains necessary.

The most concrete removal candidates are the global
\code{HasAbsUpToNormPower} propagation layer, normalization and positivity
lemmas for products used only through height, the
\code{quotientRatio}/\code{turnRatio} endpoint descriptions, and the named
\code{rallisProduct} interface in this branch. Their replacements still
require local cost proofs and a readable presentation of the triangle.
The implementation plan is to count deletions only after the old
declarations have no remaining callers.

A conservative declaration audit isolates 27 purely representational
declarations: 136 physical declaration/body lines, or 198 lines including
their attached documentation. The broader 306-line operator-product
region in \code{NormOfChain} overlaps that count and includes adapters
that require replacement. The audit therefore does not support adding
these figures or claiming a numerical net saving. Exact names and
ranges are in the archived metadata audit.

The gain is architectural: one certificate interface replaces repeated
exact scalar descriptions across seven concrete consumers and the generic
triangle. It also removes coupling between local arithmetic identities
and the final height estimate. A faithful readable implementation will
be longer than the highly compressed compact implementation; comparing
its 63-line cost file directly with a 430-line readable file is not a fair
net-saving calculation.
